% Part: first-order-logic
% Chapter: proof-systems
% Section: tableaux

\documentclass[../../../include/open-logic-section]{subfiles}

\begin{document}

\iftag{FOL}
      {\olfileid{fol}{prf}{tab}}
      {\olfileid{pl}{prf}{tab}}

\olsection{\usetoken{P}{tableau}}

Veel !!{derivation}ssystemen werken met rangschikkingen van
!!{sentence}s; !!{tableau}s werken daarentegen met !!{signed formula}s.
!!^a{signed formula} is een paar dat bestaat uit een
waarheidswaardeteken ($\True$ of $\False$) en !!a{sentence}:
\[
\sFmla{\True}{!A} \text{ of } \sFmla{\False}{!A}.
\]
!!^a{tableau} bestaat uit !!{signed formula}s die in een neerwaarts
vertakkende boom zijn gerangschikt. Het begint met een aantal
\emph{aannames} en wordt voortgezet met !!{signed formula}s die door
toepassing van een afleidingsregel uit een erboven staande
!!{signed formula} volgen. Elke regel staat toe dat aan het uiteinde
van een tak een of meer !!{signed formula}s worden toegevoegd, of dat
twee !!{signed formula}s naast elkaar worden toegevoegd. In het
laatste geval splitst de tak zich in tweeën en vormen de twee
toegevoegde !!{signed formula}s de uiteinden van de nieuwe takken.

Wanneer een regel op een samengestelde !!{signed formula} wordt
toegepast, worden !!{signed formula}s toegevoegd die onmiddellijke
deel-!!{formula}s zijn. Voor elk van beide tekens is er een regel. De
regel $\TRule{\True}{\land}$ wordt bijvoorbeeld toegepast op
$\sFmla{\True}{!A \land !B}$. Daarmee mogen beide !!{signed formula}s
$\sFmla{\True}{!A}$ en~$\sFmla{\True}{!B}$ worden toegevoegd aan het
uiteinde van elke tak die $\sFmla{\True}{!A \land !B}$ bevat. Met de
regel $\TRule{\False}{!A \land !B}$ mag een tak worden gesplitst door
$\sFmla{\False}{!A}$ en $\sFmla{\False}{!B}$ naast elkaar toe te
voegen. !!^a{tableau} is gesloten als elk van zijn takken een bij
elkaar horend paar !!{signed formula}s $\sFmla{\True}{!A}$ en
$\sFmla{\False}{!A}$ bevat.

De op !!{tableau}s gebaseerde relatie $\Proves$ wordt als volgt
gedefinieerd: $\Gamma \Proves !A$ dan en slechts dan als er een eindige
verzameling~$\Gamma_0 = \{!B_1, \dots, !B_n\} \subseteq \Gamma$ bestaat
waarvoor er een gesloten !!{tableau} is met de aannames
\[
\{\sFmla{\False}{!A}, \sFmla{\True}{!B_1}, \dots, \sFmla{\True}{!B_n}\} 
\]
De volgende gesloten !!{tableau} laat bijvoorbeeld zien dat $\Proves
(!A \land !B) \lif !A$:
\begin{oltableau}
  [\sFmla{\False}{(\formula{A} \land \formula{B}) \lif \formula{A}}, just = \TAss
    [\sFmla{\True}{\formula{A} \land \formula{B}}, just = {\TRule{\False}{\lif}[1]}
      [\sFmla{\False}{\formula{A}}, just = {\TRule{\False}{\lif}[1]}
        [\sFmla{\True}{\formula{A}}, just = {\TRule{\True}{\lif}[2]}
          [\sFmla{\True}{\formula{B}}, just = {\TRule{\True}{\lif}[2]}, close]
        ]
      ]
    ]
  ]
\end{oltableau}

Een verzameling $\Gamma$ is inconsistent in de !!{tableau}calculus als
er een gesloten !!{tableau} bestaat met aannames
\[
\{\sFmla{\True}{!B_1}, \dots, \sFmla{\True}{!B_n}\} 
\]
waarbij $!B_i \in \Gamma$ geldt voor elke gebruikte index.

!!^{tableau}s zijn in de jaren vijftig van de twintigste eeuw
onafhankelijk van elkaar uitgevonden door Evert Beth en Jaakko
Hintikka en vervolgens vereenvoudigd en verbreid door Raymond
Smullyan. Ze zijn zeer gemakkelijk te gebruiken, omdat het opstellen
van !!a{tableau} een zeer systematische procedure is. Door hun
systematische aard lenen !!{tableau}s zich ook voor implementatie op
een computer. !!^a{tableau} is echter vaak moeilijk te lezen en het
verband met bewijzen is soms moeilijk te doorzien. De benadering is
bovendien tamelijk algemeen: voor veel verschillende logica's bestaan
!!{tableau}systemen. Met !!{tableau}s kunnen ook !!{structure}s worden
gevonden die gegeven (verzamelingen) !!{sentence}s vervullen. Als de
verzameling vervulbaar is, heeft zij geen gesloten !!{tableau}; elk
!!{tableau} heeft dan dus een open tak. Uit een open tak kan de
vervullende !!{structure} worden ``afgelezen'', mits op die tak elke
toepasbare regel is toegepast. Er bestaat ook een zeer nauw verband
met de sequentencalculus: in wezen is een gesloten !!{tableau} een
samengeperste !!{derivation} in de sequentencalculus, ondersteboven
geschreven.
\end{document}
